\begin{abstract}
We study four questions on CIFAR-10~\cite{krizhevsky2009cifar} with controlled comparisons rather than
leaderboard numbers. A fully connected network fits the training set but generalises poorly; a CNN with the
same parameter count and training budget gains 32 test points, while four times more MLP parameters gain 1.4.
Training on partly random labels shows that the MLP learns transferable structure in the first few epochs
and memorises afterwards. A fixed pixel permutation leaves the MLP unchanged but costs the CNN most of its
advantage, and shuffling image patches separates local evidence from global layout. For autoencoders,
reconstruction error and class information only agree roughly, and denoising recovers part, not all, of what a
fixed classifier needs. Finally an SVM and a clustering of the same autoencoder code are compared class pair by
class pair. All numbers quoted as final results are test-set numbers obtained after every choice was fixed on
validation data.
\end{abstract}

\section{Common setup}
\label{sec:setup}
\textbf{Data and split.} Images come from the CIFAR-10 parquet release \texttt{uoft-cs/cifar10} (revision
\texttt{0b27149}); it holds the same images as the official batches (identical channel statistics, test set in
the official order) but stores the training images in a different order, to which the saved split indices
refer. The 50{,}000 training images were split once into a stratified 45{,}000 /
5{,}000 train/validation split (4{,}500/500 per class) with seed \texttt{242003}, the digits of the roll number
M24MA2003 in order (\texttt{common/split\_seed242003.npz}). Every question loads this file. Inputs are normalised
per channel with the mean and standard deviation of the 45k training split. The validation split is used for
every model, architecture and hyper-parameter choice; the official 10k test set is evaluated only after the
choices are fixed, once per selected model.

\textbf{Training protocol (shared by Q1 and Q2).} 40 epochs, batch 128, SGD with Nesterov momentum 0.9, weight
decay $5\!\times\!10^{-4}$, OneCycle schedule (15\% warm-up), seed 242003. The peak learning rate is chosen per
model from $\{0.01, 0.03, 0.1, 0.3, 1.0\}$ by final-epoch validation accuracy. The grid started as
$\{0.01,\dots,0.1\}$ and was widened twice because a model was selected at its upper edge; every model was then
run on the full five-value grid, so tuning effort is identical. Train accuracy in all curves is measured in
evaluation mode on the un-augmented training images. Each configuration is run with one seed, so differences of
one to two points should not be over-read. Implementation: PyTorch~\cite{paszke2019pytorch} and
scikit-learn~\cite{pedregosa2011scikit}; no pretrained weights or external data.

\textbf{Reuse across questions.} Table~\ref{tab:reuse} lists every model or result used outside the question
that introduced it.

\begin{table}[t]
\centering
\small
\resizebox{\columnwidth}{!}{%
\begin{tabular}{@{}lll@{}}
\toprule
Component & Introduced & Reused in \\
\midrule
Split, normalisation, training loop & Setup & all \\
MLP 3072-336-336-10 and CNN (Tab.~\ref{tab:q1a}) & Q1(a) & Q2(a)--(c) \\
Clean-condition Q2 results (no retraining) & Q1(a) & Q2(a) \\
Selected Q1(a) CNN as fixed classifier & Q1(a) & Q3(b) \\
Clean AE-1024 & Q3(a) & Q3(b),(c) \\
DAE-1024 encoder (frozen) as representation & Q3(b) & Q3(c), Q4 \\
\bottomrule
\end{tabular}}
\caption{Models and results reused across questions.}
\label{tab:reuse}
\end{table}
